
\subsubsection{6.1 Key Findings}

The results support the hypothesis that attribution methods face structural trade-offs in non-convex settings. The most unexpected finding was that TracIn and FastIF---both gradient-based methods---showed the lowest Jaccard agreement (0.0024). The gradient-based categorization appears too coarse; these different implementations navigate the non-convex landscape differently even when sharing a common paradigm.

The contrast between convex and non-convex settings is consistent with the mechanistic hypothesis: in convex logistic regression, the unique global minimum and well-defined Hessian structure ensure that all approximation methods converge toward the same target, keeping metrics coupled. In non-convex deep networks, different approximation strategies explore different regions of the loss landscape.

\subsubsection{6.2 Practical Guidance}

Based on these findings, the following guidance may be considered:

\begin{table}[htbp]
\centering
\resizebox{\columnwidth}{!}{%
\begin{tabular}{lll}
\toprule
Application & Priority Metric & Potential Method Candidates \\
\midrule
Data debugging (find top-k) & Rank preservation & IF, TRAK \\
Data valuation (accurate scores) & Magnitude fidelity & FastIF, TracIn \\
ML auditing & Application-dependent & Match to audit goal \\
\bottomrule
\end{tabular}
}
\end{table}

When stakes are high, practitioners may wish to use multiple methods and acknowledge that attribution results are method-dependent.

\subsubsection{6.3 Limitations}

Several limitations of this work are acknowledged:

\begin{enumerate}
\item \textbf{Simplified implementations and ground truth approximation.} The experiments use gradient-based proxies rather than true LOO retraining for ground truth, and consistent re-implementations of attribution methods rather than official library versions. While this design enables fair compute-normalized comparison, the observed trade-offs could potentially be artifacts of specific implementations. However, the core finding---that metrics are coupled in convex settings but decoupled in non-convex settings---reflects the underlying loss landscape geometry, not specific implementation choices. The 84\% R-squared drop from convex to non-convex settings reflects how the Hessian structure changes. Different implementations may shift the Pareto frontier's location, but the existence of trade-offs is expected to persist. Validation with official libraries remains an important direction for future work.
\end{enumerate}

\begin{enumerate}
\item \textbf{Single architecture.} Only ResNet-18 was tested. The non-convex geometry argument suggests trade-offs may generalize, but empirical validation on attention-based architectures (Transformers) is needed.
\end{enumerate}

\begin{enumerate}
\item \textbf{Dataset scale.} CIFAR-10 with 5,000 samples was used, which is far from production scale. Scaling behavior of Pareto trade-offs remains unexplored.
\end{enumerate}

\begin{enumerate}
\item \textbf{Stability metric.} A third metric, normalized stability, was described in the methodology but not extensively validated. The conclusions focus on rank preservation (rho\_r) and magnitude fidelity (rho\_m) only.
\end{enumerate}

\begin{enumerate}
\item \textbf{Method coverage.} The experiments focus on foundational paradigms (TRAK, TracIn, IF, FastIF) and do not evaluate architecture-specific methods like DataInf or MAGIC. These methods may exhibit different Pareto structures in their target domains.
\end{enumerate}

\subsubsection{6.4 Broader Implications}

These findings have implications for how practitioners approach method selection. Rather than seeking the single best attribution method, practitioners may need to first identify which quality dimension matters for their use case and select accordingly. The results also suggest that when different methods identify different influential examples, this disagreement may be a structural property of the problem rather than an indication that one method is correct and others are wrong.
